% Checkpoint35 proposed source; checkpoint34 sources remain unchanged.
\statement{A.119}{Proved}{The entire residual involution cohomology is the compact product of the odd unit groups}
Let \(R\), \(E_R=F/R=P\times D_o\), and the original negative
full-GT \(\alpha\) be as in A.116.
Let \(\Gamma\in\widehat{\mathrm{GT}}_0\) be the actual normalized
full lift in A.118, put \(\beta=\Gamma^{-1}\alpha\Gamma\), and
write \(c(x)=x^{-1}\), \(c(y)=y^{-1}\).
All three actions preserve \(R\).
The actions of \(c\) and \(\beta\) on \(E_R\) coincide, whereas
the actual \(\Psi_R\) conjugates \(c\) to \(\alpha\) there.
For each \(a\in\{c,\alpha,\beta\}\) there is a canonical pointed
homeomorphism
\[
 \partial_a:E_R^a\ \xrightarrow{\ \sim\ }\ H^1(C_2,R;a),
 \qquad h\longmapsto[w^{-1}a(w)],\quad w\in F',\ wR=h.
 \tag{A.119.1}
\]
In particular
\[
 \begin{gathered}
 E_R^c=E_R^\beta
   =\prod_{p\ {\rm odd}}
       \{\operatorname{diag}(t,t^{-1}):t\in\mathbb Z_p^\times\},\\
 E_R^\alpha=\Psi_R(E_R^c),\\
 H^1(C_2,R;a)\ \cong\ \prod_{p\ {\rm odd}}\mathbb Z_p^\times
           \qquad(a=c,\alpha,\beta)
 \end{gathered}
 \tag{A.119.2}
\]
as pointed compact spaces, with the identity-unit tuple as basepoint.
The last homeomorphism uses the specified \(\Psi_R\) for \(\alpha\).
This computation does not assert conjugacy of the three full
actions or vanishing of their residual lifting obstructions.

\begin{proof}
\emph{A.119.1. The exact three-class input, without a basis assertion.}
The relevant accepted section-counting statement is
LM4~D.34, using LM4~D.32 for order and normalization:
for a normalized negative pure peripheral involution \(a\),
\(F\rtimes_a C_2\) has exactly THREE \(F\)-conjugacy classes
of sections. This is the precise scope of R1 used here.
It does not state that three chosen involutions generate that
extension. The complete original proof obtains exhaustiveness
from the genus-zero cusp triangle and the full Brown bound;
the referee's D.34 audit retains this same distinction.

Write \(a=(-1,\phi_a)\) and let \(U_a\in F'\) be its canonical
norm solution \(U_a=\phi_a a(U_a)\).
In \(F\rtimes_a\langle t\rangle\) the three elements
\[
               t,\qquad xt,\qquad (U_a^{-1}yU_a)t
 \tag{A.119.3}
\]
are involutions. Indeed \(a(x)=x^{-1}\), and the norm identity
gives
\(a(U_a^{-1}yU_a)=U_a^{-1}y^{-1}U_a\).
Their mod-two abelian labels are \(0,e_x,e_y\), respectively,
so they are distinct and exhaust the three classes.

Let \(z\in F'\) satisfy \(z\,a(z)=1\).
The involution \(zt\) is \(F\)-conjugate to one of (A.119.3).
In the second or third case, abelianization would give
\[
          0=2g_{\rm ab}+e_x\quad\text{or}\quad
          0=2g_{\rm ab}+e_y
               \quad\text{in }\widehat{\mathbb Z}^2,
\]
which is impossible modulo two.
It is therefore conjugate to \(t\) by \(g\in F\).
Now \(0=2g_{\rm ab}\), and multiplication by two on
\(\widehat{\mathbb Z}^2\) is injective, so \(g\in F'\).
Equivalently \(z=w^{-1}a(w)\) for \(w=g^{-1}\in F'\).
We have proved the ENTIRE assertion
\[
                         H^1(C_2,F';a)=\{*\}.
 \tag{A.119.4}
\]
This uses class exhaustion, not full generation or a profinite
free-product presentation for the unknown action \(a\).

The unique norm solution in A.39 also gives \(F^a=1\).
Any fixed element has zero abelianization because \(a\) acts
as \(-I\). Multiplying \(U_a\) on the right by that fixed
element gives another \(F'\)-norm solution, so uniqueness
forces it to be one. This applies to \(\alpha\) and to its
normalized conjugate \(\beta\) in A.118; for \(c\) it also follows
from the usual finite-factor centralizer theorem.

\emph{A.119.2. Every residual cocycle and every quotient fixed point.}
For \(h\in E_R^a\), its abelianization is fixed by \(-I\)
and hence zero. Thus it has a lift \(w\in F'\).
The defect \(w^{-1}a(w)\) belongs to \(R\) and satisfies the
two-factor cocycle identity. Replacing \(w\) by \(wn\),
\(n\in R\), changes it to \(n^{-1}w^{-1}a(w)a(n)\),
the prescribed equivalence.
Conversely every \(R\)-cocycle is an \(F'\)-coboundary by
(A.119.4), so all of \(H^1(C_2,R;a)\) occurs.
If defects of \(w_1,w_2\) are equivalent through \(n\in R\),
then \(w_1nw_2^{-1}\in(F')^a=1\), proving injectivity.

The compact closed lift set
\(\{w\in F':w^{-1}a(w)\in R\}\) supplies a continuous map
onto \(E_R^a\), and its defect-to-orbit map descends continuously.
The cocycle space is compact and its orbit relation under
compact \(R\) is closed. The resulting quotient is Hausdorff.
This proves (A.119.1) as a pointed homeomorphism without
choosing continuous lifts.
In particular a quotient fixed point has an actual fixed
lift in \(F'\) exactly when it is the identity.

\emph{A.119.3. The actual standard fixed group.}
The standard \(c\)-action on \(P=F^{(2)}\) inverts its two
free generators. There is an actual isomorphism
\[
 P\rtimes\langle c\rangle
       \cong C_2(t_0)\amalg_{\mathrm{pro}\text{-}2}
                       C_2(t_1)\amalg_{\mathrm{pro}\text{-}2}C_2(t_2).
 \tag{A.119.5}
\]
The map from right to left sends \(t_0,t_1,t_2\) to \(c,xc,yc\).
Conversely freeness sends \(x,y\) to \(t_1t_0,t_2t_0\);
conjugation by \(t_0\) inverts these elements, so this extends
to the semidirect product. Both composites fix the displayed
generators.
Herfort--Ribes Theorem B\(^{\prime}\), with the class of all
finite two-groups, gives the centralizer of \(t_0\) as its
own factor.\citep{A_HR1985}
That class is closed under subgroups, quotients, and extensions,
exactly as required in their \(\S2\) and full B\(^{\prime}\)
proof. Hence \(P^c=1\).

On the odd abelian coordinates of \(D_o\), \(c\) is \(-I\),
so its fixed part is zero. On \(G_p\) it is conjugation by
\(D=\operatorname{diag}(1,-1)\), since this inverts the two
actual elementary matrices \(X,Y\).
For odd \(p\), commuting with \(D\) forces both off-diagonal
entries to vanish. The fixed subgroup in \(G_p\) is therefore
\(\{\operatorname{diag}(t,t^{-1}):t\in\mathbb Z_p^\times\}\).
At three its reduction is \(I\) or \(-I\), whose \(\chi\)-value
is zero. It thus belongs to \(G_3'\) and satisfies the actual
fiber-product constraint with the zero abelian coordinate.
All such coordinate tuples occur in \(D_o\).
This proves the first line of (A.119.2).

\emph{A.119.4. The exact twisted identifications.}
A.116 proves the actual quotient identity
\(\alpha\Psi_R=\Psi_Rc\). Consequently
\(E_R^\alpha=\Psi_R(E_R^c)\).
A.118 supplies a full lift \(\Gamma\) inducing \(\Psi_R\) and
defines \(\beta=\Gamma^{-1}\alpha\Gamma\); its quotient action
is exactly \(c\), so \(E_R^\beta=E_R^c\).
This proves the remaining homeomorphisms in (A.119.2).

Their functorial marking is explicit. For \(h\in E_R^c\)
and a lift \(w\in F'\), the \(\beta\)-boundary is
\([w^{-1}\beta(w)]\). The lift \(\Gamma(w)\) of \(\Psi_R(h)\)
has \(\alpha\)-defect
\[
       \Gamma(w)^{-1}\alpha(\Gamma(w))
                    =\Gamma\bigl(w^{-1}\beta(w)\bigr).
 \tag{A.119.6}
\]
Thus the actual conjugation of the \(\alpha,\beta\) actions
on \(R\) respects the stated parameter spaces.
The identification between the \(c\)- and \(\beta\)-spaces
is instead the two boundary homeomorphisms through \(E_R^c\).
It is not a claimed automorphism of \(R\) intertwining those
actions.

For any residual \(a\)-cocycle \(z\), its unique ambient
\(F'\)-solution \(w^{-1}a(w)=z\) gives its precise obstruction
label \(wR\). The equation has a solution in \(R\) exactly
when that label is the identity. Nonidentity odd-unit labels
therefore give nontrivial actual classes.
This does not evaluate a particular original coherent
frame class or force the canonical residual word in A.118
to vanish. That word satisfies its affine norm equation
with \(\phi_\beta\), not an asserted unweighted pentagon;
the transported pentagon retains every correction printed
in A.118. Full conjugacy and the original CM sign remain open.
\end{proof}

\begingroup\footnotesize\raggedright
Primary and original-source audit: the entire LM4~D.34 proof
and its full referee report were reread in checkpoint35, with
LM4~D.32 used only for its stated normalization/order input.
Herfort--Ribes (1985), \(\S2\), Theorem B\(^{\prime}\) and
Lemma~3 with their complete proof, were freshly read in the
opened full primary text; no general profinite
subgroup Kurosh assertion is used.\par\endgroup
